\chapter{Kỹ thuật nâng cao đã chọn: hybrid search và xếp hạng lại}
\label{ch:nang-cao}

\section{Vì sao chọn}

Câu hỏi về quy chế chứa ký hiệu và con số mà embedding nắm kém (``B+'', ``Điều 10'', ``150 tín chỉ''),
và bi-encoder chỉ so khớp ở mức ý chính. Hai kỹ thuật thường được khuyến nghị cho đúng hai điểm yếu này
là trộn BM25 (hybrid search) và xếp hạng lại bằng cross-encoder. Đồ án cài cả hai và \emph{đo} xem
chúng có giúp trên corpus của mình không.

\section{Cơ chế}
\label{sec:hybrid}

Module \texttt{src/retriever.py} tổ chức truy hồi thành ba tầng độc lập, đúng theo cách LlamaIndex
tách \emph{retriever} khỏi \emph{node postprocessor}:

\begin{enumerate}
  \item \textbf{Vector top-$k$ (mặc định)}: \texttt{index.as\_retriever(similarity\_top\_k=10)}.
        Câu hỏi được embedding bằng đúng mô hình đã dùng khi lập chỉ mục, rồi so cosine với toàn bộ
        Node (\texttt{SimpleVectorStore} tìm vét cạn, phù hợp với kho vài chục đến vài trăm Node).
  \item \textbf{Hybrid BM25 + vector (tuỳ chọn, \texttt{use\_hybrid})}: câu hỏi về quy chế chứa nhiều
        ký hiệu và con số mà embedding nắm kém (``B+'', ``Điều 10'', ``150 tín chỉ'', ``QĐ 3183'').
        Đồ án cài BM25 Okapi~\cite{robertson2009bm25} thuần Python trong \texttt{src/bm25.py} ---
        gói \texttt{BM25Retriever} chính thức nằm ở một package riêng chứ không có trong
        \texttt{llama-index-core}, nên tự viết vừa nhẹ vừa cho phép kiểm soát cách tách từ.
        Hai nguồn được trộn bằng \texttt{QueryFusionRetriever} ở chế độ \texttt{reciprocal\_rerank}
        (RRF), với \texttt{num\_queries=1} để không gọi LLM sinh truy vấn phụ và giữ phép đo tất định.
  \item \textbf{Postprocessor}: ngưỡng tương đồng \texttt{SimilarityPostprocessor} (tuỳ chọn) rồi
        tới xếp hạng lại bằng cross-encoder~\cite{nogueira2019passage}
        (\texttt{SentenceTransformerRerank}, giữ top-3).
\end{enumerate}

\paragraph{Tách từ cho BM25.} \texttt{tach\_tu()} bỏ dấu tiếng Việt trước khi chấm điểm
(\texttt{Điều} $\rightarrow$ \texttt{dieu}). Nhờ vậy câu hỏi gõ thiếu dấu hoặc lỗi OCR vẫn khớp văn
bản có dấu --- đúng loại lỗi đã gặp trong nhật ký hỏi đáp. Đánh đổi là mất phân biệt các cặp chỉ
khác dấu, hiếm gặp trong văn bản quy chế. \texttt{tests/test\_bm25.py} khóa lại hành vi này bằng ba
ca: khớp theo ký hiệu ``B+'', khớp khi câu hỏi không dấu, và trả về rỗng khi không có từ nào khớp
(tránh đưa Node nhiễu vào prompt).

\paragraph{Một chi tiết dễ sai.} Điểm sau khi trộn RRF không cùng thang với cosine, nên nếu đặt
\texttt{similarity\_cutoff} khi đang bật hybrid thì mọi Node đều bị loại. Hàm
\texttt{get\_postprocessors()} vì thế chủ động bỏ qua ngưỡng khi \texttt{use\_hybrid=True} và ghi log
cảnh báo, thay vì trả về danh sách rỗng một cách im lặng.

\paragraph{Chạy trên GPU 4\,GB và trên máy từ xa.} Reranker nạp ở fp16 (Mục~\ref{sec:phan-cung}). Khi đặt
\texttt{MODEL\_SERVER\_URL}, embedding và reranker được gọi qua HTTP tới \texttt{src/model\_server.py}
(\texttt{RemoteEmbedding}, \texttt{RemoteRerank} trong \texttt{src/remote\_models.py}); client kiểm tra server
đang chạy đúng mô hình của profile để vector không lệch với chỉ mục.

\section{Số liệu}

\begin{table}[H]
  \centering
  \small
  \begin{tabularx}{\textwidth}{l l *{4}{>{\centering\arraybackslash}X} >{\centering\arraybackslash}X}
    \toprule
    \textbf{Bộ} & \textbf{Cấu hình} & \textbf{R@1 (CI 95\%)} & \textbf{R@3} & \textbf{MRR} & \textbf{nDCG@1} & \textbf{ms/câu} \\
    \midrule
    01 (31 câu) & vector (mặc định)          & 81\% [68--94] & 100\% & 0,892 & 0,806 & 165 \\
                & + rerank                    & 81\% [65--94] & 100\% & 0,892 & 0,806 & 8\,487 \\
                & hybrid (RRF, trọng số bằng) & 77\% [61--90] & 94\%  & 0,866 & 0,774 & 155 \\
                & hybrid + rerank             & 81\% [65--94] & 100\% & 0,892 & 0,806 & 8\,193 \\
                & hybrid ưu tiên vector 3:1   & 84\% [71--97] & 100\% & 0,914 & 0,839 & 148 \\
    \midrule
    02 (10 câu) & vector (mặc định)          & 90\% [70--100] & 90\%  & 0,910 & 0,900 & 149 \\
                & + rerank                    & 60\% [30--90]  & 100\% & 0,783 & 0,600 & 7\,854 \\
                & hybrid (RRF, trọng số bằng) & 60\% [30--90]  & 90\%  & 0,717 & 0,600 & 166 \\
                & hybrid + rerank             & 60\% [30--90]  & 90\%  & 0,750 & 0,600 & 8\,811 \\
                & hybrid ưu tiên vector 3:1   & 80\% [50--100] & 90\%  & 0,850 & 0,800 & 137 \\
    \bottomrule
  \end{tabularx}
  \caption{Truy hồi trên corpus 6 văn bản / 54 Node, profile \texttt{server}
  (\texttt{AITeamVN/Vietnamese\_Embedding}), đo trên CPU. Sinh lại bằng \texttt{make eval}.}
  \label{tab:recall}
\end{table}

\paragraph{2. Xếp hạng lại (rerank) không cải thiện trên kho nhỏ, và có hại cho bộ 02.}
Trên bộ 01, rerank giữ nguyên cả ba chỉ số (81\% / 100\% / 0,892) nhưng đổi thứ tự top-1 ở bốn câu
(số câu lấy đúng Quy chế 2021 tăng từ 23 lên 25). Trên bộ 02, rerank \emph{làm R@1 tụt từ 90\% xuống
60\%} và MRR từ 0,910 xuống 0,783, đổi lại R@3 tăng từ 90\% lên 100\%. Nguyên nhân hợp lý: các mục
sửa đổi của QĐ 3183 rất ngắn và dùng chung từ vựng (``thang điểm'', ``điểm chữ''), nên cross-encoder
đẩy mục sửa đổi ``na ná'' lên trên mục đúng; đổi lại nó kéo được Node đúng vào top-3. Về chi phí,
trên CPU rerank làm mỗi câu từ 0,15\,s lên 8,5\,s --- chậm hơn khoảng 50 lần --- mà không đem lại chỉ
số nào tốt hơn. Kết luận thực nghiệm: \textbf{giữ \texttt{use\_rerank=False}} trên quy mô 54 Node, và
chỉ đo lại khi corpus lớn hơn nhiều (hàng trăm đến hàng nghìn Node).

\paragraph{3. Hybrid search: bản thân BM25 ở đây yếu, và cách trộn còn quan trọng hơn việc có trộn hay không.}
Cấu hình hybrid mặc định (RRF, hai nguồn trọng số bằng nhau) \emph{làm giảm} chất lượng ở cả hai bộ:
bộ 01 R@1 81\% $\to$ 77\%, R@3 100\% $\to$ 94\%; bộ 02 R@1 90\% $\to$ 60\%. Vì BM25 một mình chỉ đạt
58\%/40\% ở R@1, việc cho nó trọng số ngang với vector kéo các Node khớp từ khóa nhưng sai ngữ nghĩa
lên đầu. Để tách nguyên nhân, đồ án đo thêm biến thể \texttt{relative\_score} với trọng số ưu tiên
vector 3:1: kết quả \emph{tốt nhất trên bộ 01} (R@1 84\%, MRR 0,914, nDCG@1 0,839) nhưng lại
\emph{kém hơn vector} trên bộ 02 (R@1 80\% so với 90\%). Nói cách khác, chênh lệch giữa các biến thể
hybrid nằm trong khoảng nhiễu của bộ câu hỏi, còn tham số trộn thì quyết định kết quả nhiều hơn cả
việc ``có BM25 hay không''. Đây là lý do \texttt{use\_hybrid} mặc định vẫn là \texttt{False}: muốn
bật phải chọn fusion mode và trọng số bằng đo trên chính corpus của mình.

\paragraph{4. Kiểm định thống kê: bộ câu hỏi hiện tại chưa đủ mạnh để khẳng định các chênh lệch nhỏ.}
Kiểm định McNemar chính xác trên R@1:
\begin{itemize}
  \item vector vs BM25, bộ 01: chỉ vector đúng 9 câu, chỉ BM25 đúng 2 câu, $p = 0{,}065$;
  \item vector vs BM25, bộ 02: 5 câu so với 0 câu, $p = 0{,}062$;
  \item vector vs hybrid ưu tiên vector, bộ 01: 0 câu so với 1 câu, $p = 1{,}000$.
\end{itemize}
Như vậy khoảng cách vector--BM25 rất lớn về độ lớn nhưng \emph{chưa đạt} ngưỡng $p<0{,}05$ với 31
câu, còn chênh lệch +3 điểm phần trăm của biến thể hybrid tốt nhất hoàn toàn có thể là may mắn. Đây
là hạn chế về \emph{độ mạnh thống kê} của bộ câu hỏi, không phải về hệ thống, và là lý do phải báo
cáo khoảng tin cậy thay vì chỉ một con số.


\section{Dùng ở đâu trong đồ án}

Cả hai đều có công tắc (\texttt{use\_hybrid}, \texttt{use\_rerank}; \texttt{--rerank}, \texttt{--hybrid} trên dòng
lệnh; ô chọn trong giao diện Gradio) và \textbf{mặc định tắt} theo kết quả đo. \texttt{make eval} đo lại cả
bốn cấu hình khi corpus thay đổi.
